% =============================================================================
% Partie 1 — Description du projet
% =============================================================================
\chapter{Description du projet}

\epigraphe{Une banque ne gagne pas d'argent en prêtant : elle en gagne quand on la rembourse.}

Avant d'écrire la moindre ligne de SQL, il faut comprendre \textbf{à quoi sert} le système que tu vas
construire, \textbf{pour qui}, et \textbf{pourquoi} il doit être aussi rigoureux. Cette partie pose le
décor : le métier du risque de crédit, le rôle de la base de données dans une banque, et la place
des modèles de Machine Learning.

% -----------------------------------------------------------------------------
\section{Le titre du projet}

\begin{situation}{Mise en situation --- Ton premier jour dans l'équipe Risque}
Tu commences ton stage dans une banque en ligne qui accorde des prêts à la consommation. Aujourd'hui,
tout est éparpillé : les dossiers de prêt dans un tableur, les paiements dans un export comptable, et
les scores du modèle de Machine Learning dans le notebook d'une data scientist.

Ta responsable te dit : \emph{« Le régulateur nous demande de justifier chaque provision, le comité
des risques veut savoir chaque matin combien de prêts sont en retard, et personne ne sait quelle
version du modèle a décidé de quel prêt. Il nous faut une vraie base de données. »}

Ta mission : construire cette base, de la première table jusqu'à la détection automatique du défaut.
\end{situation}

Le titre officiel de ton projet est :
\begin{center}
\textbf{\large Un système backend SQL d'entreprise orienté Risk Analytics,\\Scoring de Crédit et MLOps
(Fintech / Banque)}
\end{center}

Chaque mot de ce titre correspond à une exigence précise :

\begin{center}\small
\begin{tabularx}{\linewidth}{@{}B{3.6cm} L@{}}
\toprule
Système backend SQL & la base PostgreSQL et toute sa logique : tables, contraintes, vues, fonctions,
triggers, procédures. Pas d'interface graphique : c'est le moteur que d'autres applications utilisent. \\
d'entreprise & des données qui ne peuvent pas être fausses (contraintes), des opérations « tout ou
rien » (transactions), une trace de chaque changement (audit), des droits d'accès (sécurité) et des
temps de réponse corrects sur des millions de lignes (index). \\
Risk Analytics & mesurer le risque du portefeuille : retards, prêts en défaut, pertes attendues,
provisions. \\
Scoring de crédit & prédire, avant d'accorder un prêt, la probabilité que l'emprunteur ne rembourse pas. \\
MLOps & faire vivre les modèles en production : savoir quelle version a décidé quoi, conserver chaque
prédiction, comparer le modèle en place à ses concurrents, détecter qu'il se dégrade. \\
Fintech / Banque & un secteur réglementé : chaque chiffre doit pouvoir être justifié devant un
auditeur ou un superviseur. \\
\bottomrule
\end{tabularx}
\end{center}

% -----------------------------------------------------------------------------
\section{Contexte et enjeux métier}

\subsection{Le métier : prêter de l'argent qu'on n'a pas encore récupéré}

\begin{imagine}
Tu prêtes 100~€ à chacun de tes 10 amis. Tu sais d'expérience qu'en moyenne \textbf{un sur dix} ne te
rembourse pas. Et quand quelqu'un ne rembourse pas, tu finis quand même par récupérer un peu plus de la
moitié de la somme : tu perds \textbf{45~\%} de ce qu'il te doit.

\textbf{Combien t'attends-tu à perdre ?} Sur 10 amis, environ 1 ne rembourse pas. Il te devait 100~€, tu
en perds 45~\%. Ta perte attendue est donc d'environ $10 \times 10~\% \times 45~\% \times 100~€ = 45~€$.

Tu viens de faire, sans le savoir, le calcul central du risque de crédit.
\end{imagine}

Une banque fait exactement la même chose, mais avec des milliers de clients, des prêts qui durent des
années et une obligation légale de mesurer ce risque correctement. Trois grandeurs suffisent à le
décrire.

\begin{definition}{Les trois paramètres du risque de crédit}
\begin{itemize}
\item \textbf{PD} (\emph{Probability of Default}, probabilité de défaut) : la probabilité que
l'emprunteur fasse défaut. C'est ce que prédit ton modèle de scoring.
\item \textbf{LGD} (\emph{Loss Given Default}, perte en cas de défaut) : la part de la somme due que la
banque perd réellement si le client fait défaut, après recouvrement. Dans ce projet : \textbf{45~\%}, la
valeur réglementaire de référence pour une créance non garantie.
\item \textbf{EAD} (\emph{Exposure At Default}, exposition au moment du défaut) : la somme que le client
doit encore au moment où il fait défaut. Pour un prêt amortissable, c'est le \textbf{capital restant
dû} : il baisse à chaque mensualité payée.
\end{itemize}
\end{definition}

Leur produit donne la \textbf{perte attendue} (\emph{Expected Loss}) :
\[
\boxed{\;EL = PD \times LGD \times EAD\;}
\]

\begin{calcul}{la perte attendue d'un prêt}
Un client emprunte 10\,000~€. Ton modèle estime sa probabilité de défaut à 5~\%. La LGD vaut 45~\%.
\begin{enumerate}
\item EAD : au départ, le client doit tout le capital, soit \textbf{10\,000~€}.
\item Perte si défaut : $45~\% \times 10\,000 = 4\,500$~€.
\item Perte attendue : $5~\% \times 4\,500 = \mathbf{225~€}$.
\end{enumerate}
Ces 225~€ ne sont pas une perte certaine : c'est la perte \textbf{moyenne} sur un grand nombre de prêts
semblables. La banque doit la couvrir par le taux d'intérêt qu'elle facture et par des provisions.
\end{calcul}

\subsection{Les trois questions de l'équipe Risque}

Tout le travail d'une équipe Risque de crédit répond à trois questions, posées à trois moments de la vie
d'un prêt. Ton projet apporte une réponse SQL à chacune.

\begin{center}\small
\begin{tabularx}{\linewidth}{@{}B{3.9cm} P{2.6cm} L@{}}
\toprule
Question & Moment & Réponse dans ton projet \\
\midrule
Qui \emph{va} faire défaut ? & avant l'octroi & le \textbf{score} du modèle (PD), stocké dans
\code{predictions_scoring_ml} avec son explication SHAP. \\
Qui \emph{est en train} de faire défaut ? & pendant la vie du prêt & le \textbf{retard de paiement} (DPD),
calculé à partir de \code{echeancier_remboursements}. \\
Combien allons-nous perdre ? & chaque fin de mois & l'\textbf{exposition} (EAD) et les
\textbf{provisions} (ECL), calculées et historisées dans \code{provisions_mensuelles}. \\
\bottomrule
\end{tabularx}
\end{center}

\subsection{Le cadre réglementaire, en deux idées}

Tu n'as pas besoin d'être juriste, mais deux textes structurent tout le projet. Les citer correctement en
entretien fait une vraie différence.

\begin{definition}{Le défaut (règlement européen CRR, article 178)}
Un emprunteur est \textbf{en défaut} lorsqu'il a un retard de paiement de \textbf{plus de 90 jours} sur
une obligation significative, ou lorsque la banque estime qu'il ne paiera probablement pas
(\emph{unlikeliness to pay}). C'est la règle que ton trigger appliquera automatiquement.
\end{definition}

\begin{definition}{Les provisions IFRS 9}
Depuis 2018, la norme comptable IFRS~9 oblige les banques à provisionner les \textbf{pertes attendues}
(\emph{Expected Credit Loss}, ECL) \textbf{dès l'octroi} du prêt, sans attendre l'incident. Chaque prêt
est classé dans l'un de trois \textbf{stages} :
\begin{itemize}
\item \textbf{Stage 1} : prêt sain. On provisionne la perte attendue sur les \textbf{12 prochains mois}.
\item \textbf{Stage 2} : risque significativement dégradé, présumé à partir de \textbf{plus de 30 jours}
de retard. On provisionne la perte attendue sur \textbf{toute la durée restante}.
\item \textbf{Stage 3} : prêt en défaut. La perte est considérée comme avérée : $PD = 100~\%$.
\end{itemize}
\end{definition}

\begin{observation}[Pourquoi c'est un sujet SQL]
Le superviseur et les auditeurs peuvent demander, des mois plus tard : \emph{« Pourquoi avez-vous
provisionné 57,21~€ sur le prêt 61 au 30 septembre ? »}. Il faut alors retrouver la donnée utilisée,
la règle appliquée et la date du calcul. Un tableur modifié à la main en est incapable. Une base de
données bien conçue, avec ses contraintes, son historique et ses calculs rejouables, le permet.
\end{observation}

\subsection{Où intervient le Machine Learning}

Ton projet de scoring a produit trois modèles, entraînés en R sur l'historique des prêts : une
\textbf{grille de score} (des points par caractéristique), une \textbf{régression logistique}, et un
\textbf{XGBoost monotone}. Sur l'échantillon de test, le XGBoost atteint une AUC de \textbf{0,923},
contre 0,876 pour la logistique : c'est lui qui décide en production. On l'appelle le
\textbf{modèle champion} ; les deux autres sont des \textbf{challengers}.

\begin{definition}{SHAP}
Les valeurs SHAP (\emph{SHapley Additive exPlanations}) décomposent une prédiction en
\textbf{contributions de chaque variable}. Pour un prêt donné, une valeur positive pousse le risque
vers le haut, une valeur négative le fait baisser. Ta base conserve, pour chaque prédiction du
champion, la variable qui a le plus pesé et sa contribution.
\end{definition}

Par exemple, pour le prêt 61, le XGBoost prédit une PD de 1,07~\% et la variable la plus influente est
la \code{note_risque}, avec une contribution de $-1{,}29$ : sa bonne note (A) a fortement fait baisser
le risque estimé. Cette explication est indispensable : le RGPD encadre les décisions automatisées, et
les orientations de l'Autorité bancaire européenne sur l'octroi de prêts exigent des modèles
explicables. Une banque doit pouvoir dire à un client \emph{pourquoi} son prêt a été refusé.

La base de données est le point de rencontre entre le monde des modèles et celui du métier. La
figure~\ref{fig:boucle-mlops} montre ce cycle, que l'on appelle souvent la \textbf{boucle MLOps}.

\begin{figure}[ht]
\centering
\begin{tikzpicture}[font=\small, >=Stealth,
  etape/.style={draw=bleuML, fill=bleuFond, rounded corners=3pt, minimum width=3.3cm, minimum height=1cm,
                align=center, line width=0.6pt}]
  \node[cylinder, shape border rotate=90, aspect=0.2, draw=bleuML, fill=white, line width=0.9pt,
        minimum width=2.6cm, minimum height=1.7cm, align=center] (db) at (0,0) {\textbf{Base}\\\textbf{PostgreSQL}};
  \node[etape] (a) at (90:3.1)  {\textbf{1.} Données d'historique\\{\footnotesize clients, prêts, défauts}};
  \node[etape] (b) at (30:3.6)  {\textbf{2.} Entraînement\\{\footnotesize R / Python, XGBoost}};
  \node[etape] (c) at (-30:3.6) {\textbf{3.} Prédiction\\{\footnotesize PD + explication SHAP}};
  \node[etape] (d) at (-90:3.1) {\textbf{4.} Décision et suivi\\{\footnotesize échéancier, paiements}};
  \node[etape] (e) at (-150:3.6){\textbf{5.} Vérité observée\\{\footnotesize défaut ou non}};
  \node[etape] (f) at (150:3.6) {\textbf{6.} Monitoring\\{\footnotesize AUC, déciles, champion}};
  \draw[->, thick, bleuML] (a) to[bend left=12] (b);
  \draw[->, thick, bleuML] (b) to[bend left=12] (c);
  \draw[->, thick, bleuML] (c) to[bend left=12] (d);
  \draw[->, thick, bleuML] (d) to[bend left=12] (e);
  \draw[->, thick, bleuML] (e) to[bend left=12] (f);
  \draw[->, thick, bleuML] (f) to[bend left=12] (a);
  \foreach \n in {a,c,d,e,f} \draw[densely dotted, grisTitre] (db) -- (\n);
\end{tikzpicture}
\caption{La boucle MLOps : la base fournit les données d'entraînement, conserve chaque prédiction,
enregistre ce qui se passe réellement, puis permet de mesurer la performance du modèle.}
\label{fig:boucle-mlops}
\end{figure}

\begin{analogie}{le carnet de santé}
Pense au prêt comme à un patient. Le \textbf{modèle de scoring} est le médecin qui fait un pronostic à
l'entrée. La \textbf{base de données} est le carnet de santé : on y note le pronostic, qui l'a fait et
quand, puis chaque examen (chaque mensualité). Le \textbf{suivi du retard} est la prise de température
régulière, et le \textbf{trigger} est l'alarme qui se déclenche quand la fièvre dépasse un seuil. Sans
carnet, impossible de savoir, des années plus tard, si le médecin avait vu juste.
\end{analogie}

\subsection{Pourquoi ce projet existe en entreprise}

\begin{center}\small
\begin{tabularx}{\linewidth}{@{}B{3.1cm} L L@{}}
\toprule
Besoin & Avec des fichiers et un notebook & Avec ton système SQL \\
\midrule
Données fiables & un âge de 15 ans ou un paiement en double passent inaperçus & les contraintes les
refusent à l'entrée \\
Source unique & chacun a « sa » version du fichier & une seule base, partagée \\
Traçabilité & on ne sait plus quel modèle a décidé & chaque prédiction est conservée avec sa version \\
Détection du défaut & quelqu'un doit y penser & un trigger et un traitement de nuit l'appliquent \\
Provisions & recalculées à la main, non reproductibles & une procédure rejouable, historisée chaque mois \\
Volume & le tableur sature vers un million de lignes & 1,1 million de mensualités interrogées en
quelques millisecondes grâce aux index \\
Sécurité & le fichier circule par e-mail & des rôles : l'analyste lit, l'application encaisse \\
\bottomrule
\end{tabularx}
\end{center}

% -----------------------------------------------------------------------------
\section{Cas d'usage concret : la vie d'un prêt}

\begin{situation}{Le dossier d'Awa Koné}
Awa Koné, 29 ans, locataire, 42\,000~€ de revenus annuels, demande \textbf{12\,000~€ sur 36 mois} au
taux de 9,5~\% pour financer ses études. Suivons son prêt de la demande jusqu'au défaut éventuel, et
regardons ce que fait la base à chaque étape.
\end{situation}

\begin{enumerate}[leftmargin=1.6em, itemsep=0.45em]
\item \textbf{La demande (20/09/2026).} Une fiche client et un dossier de prêt sont créés. Les contraintes
vérifient tout de suite que l'âge est plausible, que le montant est positif, que le motif fait partie
de la liste autorisée.
\item \textbf{Le scoring (24/09/2026).} Le XGBoost estime sa probabilité de défaut à \textbf{11,8~\%}. Le
seuil de décision du modèle est 25,5~\% : le prêt est \textbf{accordé}. L'explication SHAP indique que
la variable la plus influente est le \textbf{taux d'effort} (le montant emprunté rapporté au revenu,
ici $12\,000 / 42\,000 \approx 29~\%$), qui augmente le risque. La prédiction est enregistrée avec la
version du modèle, le seuil et la date.
\item \textbf{L'octroi et l'échéancier (25/09/2026).} La base génère \textbf{36 mensualités de
384,40~€} à partir du 25/10/2026. Le coût total du crédit est de 1\,838,21~€ d'intérêts.
\item \textbf{Le suivi.} Chaque paiement est enregistré par une procédure qui l'affecte à la mensualité la
plus ancienne restant due. Le retard (DPD) se déduit des dates : il n'est jamais saisi à la main.
\item \textbf{Le retard.} Si la mensualité du 25/10/2026 n'est pas payée, Awa est \textbf{en retard} dès
le lendemain. Au-delà de 30 jours, son prêt passe en \textbf{stage 2} et la provision augmente.
\item \textbf{Le défaut.} Toujours sans paiement, au 31/01/2027 son retard atteint \textbf{98 jours} :
plus de 90 jours, le prêt passe \textbf{en défaut}. Le changement est inscrit dans un journal d'audit,
une alerte est envoyée, et la provision est calculée avec une PD de 100~\%.
\item \textbf{La fin.} Si tout est remboursé, le prêt est \textbf{soldé}. S'il est en défaut, il ne
redevient pas sain automatiquement : une période probatoire et une décision humaine sont nécessaires.
\end{enumerate}

Ces changements d'état forment le cycle de vie du prêt (figure~\ref{fig:cycle-vie}). Tu
l'implémenteras au chapitre~3.

\begin{figure}[ht]
\centering
\begin{tikzpicture}[font=\small, >=Stealth, node distance=1.2cm,
  etat/.style={draw, rounded corners=8pt, minimum width=2.25cm, minimum height=0.85cm, line width=0.7pt, align=center}]
  \node[etat, draw=grisTitre, fill=grisFond] (att) at (0,0) {en attente};
  \node[etat, draw=grisTitre, fill=grisFond] (ref) at (0,-2.6) {refusé};
  \node[etat, draw=vertCadre, fill=vertFond] (sain) at (3.6,0) {sain};
  \node[etat, draw=orangeEx, fill=orangeFond] (ret) at (7.6,0) {en retard};
  \node[etat, draw=rouge, fill=rougeFond] (def) at (12.1,0) {en défaut};
  \node[etat, draw=bleuML, fill=bleuFond] (sol) at (7.6,-2.6) {soldé};
  \draw[->, thick] (att) -- node[left, font=\footnotesize, align=right] {PD $\geq$ seuil} (ref);
  \draw[->, thick] (att) -- node[above, font=\footnotesize] {accordé} (sain);
  \draw[->, thick] (sain) to[bend left=28] node[above=1pt, font=\footnotesize] {DPD $> 0$} (ret);
  \draw[->, thick] (ret) to[bend left=28] node[below=1pt, font=\footnotesize] {régularisation} (sain);
  \draw[->, thick, rouge] (ret) -- node[above, font=\footnotesize] {DPD $> 90$} (def);
  \draw[->, thick, dashed, grisTexte] (def) to[bend right=40] node[above=1pt, font=\footnotesize, align=center]
      {période probatoire + décision humaine} (sain);
  \draw[->, thick] (sain) |- node[pos=0.75, below, font=\footnotesize] {reste dû $= 0$} (sol);
  \draw[->, thick] (ret) -- node[right, font=\footnotesize] {reste dû $= 0$} (sol);
  \draw[->, thick, dashed, grisTexte] (def) |- node[pos=0.75, below, font=\footnotesize] {recouvrement} (sol);
\end{tikzpicture}
\caption{Le cycle de vie d'un prêt. Les flèches pleines sont automatiques (trigger ou traitement de
nuit) ; les flèches en pointillés demandent une décision humaine.}
\label{fig:cycle-vie}
\end{figure}

Le tableau suivant relie chaque étape aux tables et aux mécanismes SQL que tu vas construire.

\begin{center}\small
\begin{tabularx}{\linewidth}{@{}B{2.6cm} P{4.3cm} L P{1.5cm}@{}}
\toprule
Étape & Tables & Mécanisme SQL & Où \\
\midrule
Demande & \code{clients}, \code{demandes_pret} & contraintes \code{CHECK}, clés, transaction & Ch.~1 \\
Scoring & \code{predictions_scoring_ml} & registre versionné, contrainte de cohérence & Ch.~1 \\
Échéancier & \code{echeancier_remboursements} & \code{generate_series}, CTE, formule d'annuité & Ch.~1 \\
Suivi du retard & vue \code{v_situation_prets} & agrégats, \code{FILTER}, fonctions de dates & Ch.~2 \\
Provisions & \code{provisions_mensuelles} & CTE, \code{ROLLUP}, procédure rejouable & Ch.~2 et 3 \\
Paiement & \code{echeancier_remboursements} & procédure stockée, verrou & Ch.~3 \\
Défaut & \code{demandes_pret}, \code{journal_statuts} & triggers, traitement de nuit, \code{pg_notify} & Ch.~3 \\
Monitoring & \code{predictions_scoring_ml}, \code{echantillons_modele} & \code{NTILE}, rangs, AUC en SQL & Ch.~2 \\
\bottomrule
\end{tabularx}
\end{center}

\begin{resume}
\begin{enumerate}
\item Le risque de crédit se résume par $EL = PD \times LGD \times EAD$ : probabilité de défaut, perte
en cas de défaut, exposition.
\item L'équipe Risque répond à trois questions : qui \emph{va} faire défaut (scoring), qui \emph{est en
train} de faire défaut (retards), combien allons-nous perdre (provisions).
\item Deux règles structurent le projet : le \textbf{défaut} au-delà de 90 jours de retard (CRR,
art.~178) et les \textbf{trois stages} de provisionnement IFRS~9 (seuil de 30 jours pour le stage 2).
\item La base de données est le point de rencontre entre les modèles de ML et le métier : elle fournit
les données, conserve chaque prédiction et permet de mesurer la performance des modèles.
\item Le cas d'usage suit un prêt de la demande au défaut : chaque étape correspond à une table et à un
mécanisme SQL du cours.
\end{enumerate}
\end{resume}
